\documentclass[11pt,a4paper]{article}

\usepackage[T1]{fontenc}
\usepackage[utf8]{inputenc}
\usepackage{amsmath,amssymb}
\usepackage{booktabs}
\usepackage{listings}
\usepackage{url}
\usepackage[margin=1in]{geometry}
\usepackage[hidelinks]{hyperref}

\lstset{basicstyle=\ttfamily\small,columns=fullflexible,breaklines=true}

\newcommand{\code}[1]{\texttt{#1}}

\title{\textbf{libgmpecpp}: An Elliptic Curve Primality Prover as a
  Reusable Library, with a CGBN/CUDA Backend for the Factoring Stage}

\author{Stefan Teleman \code{<stefan.teleman@gmail.com>}\\[4pt]Claude Opus 5.5 (Anthropic)}

\date{September 3, 2026; revised November 17, 2026}

\begin{document}
\maketitle

\begin{abstract}
We describe the design, implementation and empirical evaluation of
\emph{libgmpecpp}, a shared library providing elliptic curve primality
proving (ECPP) in the Atkin--Morain formulation, together with a GPU
backend for its factoring stage built on NVIDIA's CGBN fixed-width
multiple-precision library.

The starting point is a single C++ translation unit descended from
Pate Williams' 1997 implementation of Cohen's Algorithm 9.2.4, with
later contributions by James Wanless, Chris Studholme, Elisavet
Konstantinou, Yiannis Stamatiu, Christos Zaroliagis, Scott Contini, Paul
Zimmermann and Timo Poikola. Converting that program into a library
required three changes that are individually small but jointly
structural: the removal of all direct console output in favour of a
line-buffered logging callback; the elimination of a process-terminating
\code{exit(2)} deep inside the polynomial arithmetic, which we replace by
an exception carrying the nontrivial divisor that a failed modular
inversion exhibits; and the return of the primality certificate as
structured data rather than as printed text, which additionally permits
the library to distinguish outcomes that the historical command line
interface conflates under a single exit status. We present a C API,
stable across C and C++, now exporting thirty-six symbols. The original
command line tool is reimplemented as a client of this API and is shown
to be byte-identical to the historical program across 156 runs spanning
thirteen flag combinations, four inputs and three seeds, comparing
standard output, standard error and exit status.

We then profile the prover by phase and report a result that governs the
remainder of the work: the distribution of runtime between the
class-polynomial layer and the factoring layer is not fixed, but is
controlled almost entirely by the floating-point working precision. At
the historical default of $10^4$ bits, construction of Hilbert class
polynomials accounts for $66$--$89\%$ of runtime on inputs of $30$ to
$60$ decimal digits and elliptic curve factorisation for $1$--$7\%$; at
$10^3$ bits the same measurements invert, with factorisation accounting
for $58$--$99.5\%$. The absolute cost of factorisation is unchanged
between the two settings, since it is integer work; what changes is the
cost of everything surrounding it. We reduce the default precision to
$10^3$ bits after verifying across the thirty-three primes of the committed
benchmark set, spanning $20$ to $180$ decimal digits, that every input still proves, that no
loss-of-precision event is triggered, and that no proof becomes slower.
This alone yields speedups of up to $48.5\times$ on small inputs.

Because CGBN provides fixed-width \emph{integer} arithmetic while the
class-polynomial layer is arbitrary-precision \emph{floating point}
implemented as hand-rolled Taylor series, only the factoring stage admits
a CGBN formulation. Under the revised precision default this is
fortunate rather than limiting, since factorisation is now the dominant
cost. We implement stage-one ECM on Montgomery curves under Suyama's
parametrisation, in $XZ$ coordinates so that the inner loop requires no
modular inversion, with one CGBN instance per curve and Montgomery
representation throughout. We show the resulting kernel to be latency
bound rather than throughput bound: sixteen-fold increases in
concurrent curve count cost only $1.4\times$ in wall time for stage one, since the
binding constraint is the serial chain of prime powers each curve must
traverse. Consequently curves are nearly free, which is precisely the
shape that ECM rewards.

Measured on an RTX PRO 4500 Blackwell with stage two on the GPU, whole-proof
speedups are $0.99$--$1.02\times$ below sixty decimal digits, where the host
already splits cofactors in milliseconds, rising to $2.53\times$ at ninety
digits, $3.27\times$ at one hundred and ten, and $25.6\times$ at one hundred
and forty. Across a thirty-six-proof sweep from ninety to one hundred and
eighty digits the GPU backend is faster in twenty-seven runs, with median
speedup $2.37\times$ and aggregate $3.48\times$. We argue that the aggregate
figure understates the effect, which is better characterised as variance
reduction: host runtimes span $0.83$ to $67.3$ seconds while
GPU-assisted runtimes span $0.91$ to $19.7$ seconds. The mechanism is that a discriminant the host
abandons for want of a factorisation is one the GPU can use, so
GPU-assisted proofs are frequently \emph{shorter} as well as faster.

We report two negative results that shaped the design. First, dispatching
to the GPU before the host has attempted a cofactor makes proofs
approximately twice as slow, because the majority of factoring calls in
an ECPP downrun exist in order to reject candidate curve orders quickly
and are expected to fail; the backend therefore attempts the host path
first and escalates only above a configurable cofactor width. Second,
byte-comparison against the historical program is an inappropriate
correctness criterion once the GPU is engaged, since it discovers
different factorisations and hence produces different, equally valid
certificates. We therefore supply an independent certificate verifier
sharing no code with the prover, which re-derives chain linkage, the
divisibility $q \mid m$, the Goldwasser--Kilian bound $q > (\sqrt[4]{n}+1)^2$,
Hasse's bound on the curve order, curve membership of the base point, and
the two scalar multiplication conditions, and which is itself subjected
to a negative control. Twenty certificates produced by both backends
across ten large primes verify successfully.

Finally we examine the prover for shared- and distributed-memory
parallelism. The downrun is sequential by construction, but two finer
axes are available and suit different tools. Elliptic curve trials are
independent and one success ends the search, which is shared-memory work
with a negligible critical section; threading them yields $5.98\times$
on thirty-two cores. Discriminants are likewise independent, and a
speculative parallel sweep -- kept deterministic by seeding each attempt
from its own discriminant, and kept useful by retaining the smallest
successful discriminant rather than the first to arrive -- yields
$1.45\times$ on the same inputs. Across proofs, whole numbers are
embarrassingly parallel, which suits MPI; we show that neither tool
suffices alone, since MPI plateaus at $2.63\times$ on a batch bounded by
its slowest single proof while threads reduce that proof, the hybrid
reaching $11.0\times$.

A November 2026 revision adds a second GPU entry point,
\code{gmpecpp\_cuda\_ecm()}, for programs that factor rather than prove: no host attempt and no
global lock, one full wave of curves, short resumable launches that a caller can stop and that
remove any cap on $B_1$, and a stage two on the GPU (prime continuation with normalized baby
points, two modular multiplications per prime), which raises the per-curve success rate about
tenfold at modest cost. Used beside GMP-ECM on the host, one GPU then finds about nine times as
many 30-digit factors per second as thirty-one CPU threads. A companion library,
\code{libgmpecmfft}, carries ECM past CGBN's $16352$-bit limit with cuFFT-based Montgomery
arithmetic and Montgomery's FFT continuation for stage two, validated against CGBN, GMP and
GMP-ECM on identical curves; \code{gmpecmfft-mpi} shares one number's curves among the GPUs of
a cluster.

The complete system is licensed GPL-3.0-or-later, comprises a shared
library, a compatible command line tool, an MPI driver, benchmark
harnesses and an independent verifier, and builds with or without CUDA
and MPI present.
\end{abstract}

\tableofcontents

\section{Introduction}

A probabilistic primality test such as Miller--Rabin \cite{miller1976,
rabin1980} establishes compositeness conclusively but primality only with
high probability. Elliptic curve primality proving, introduced by
Goldwasser and Kilian \cite{goldwasser1986} and made practical by Atkin
and Morain \cite{atkin1993}, instead produces a \emph{certificate}: a
chain of elliptic curves whose validity can be checked independently of
the program that produced it, and in substantially less time than it took
to construct.

This paper documents the conversion of an existing ECPP implementation
into a reusable library, the addition of a GPU backend for its factoring
stage, and the empirical characterisation of both. We emphasise the
measurements throughout, because two of the design decisions that matter
most were forced by data that contradicted our initial expectations.

\section{Provenance and attribution}

The prover descends from a program written by Pate Williams in 1997,
implementing Algorithm 9.2.4 of Cohen \cite{cohen1993}, with subsequent
contributions from James Wanless (2008--2019), Chris Studholme (2004),
Elisavet Konstantinou, Yiannis Stamatiu and Christos Zaroliagis (2003),
Scott Contini (1996), Paul Zimmermann (2005) and Timo Poikola (2010). It
is distributed under the GNU General Public License, version 3 or later,
and the present work preserves that licensing.

The mathematical content is almost entirely classical or due to the
authors cited in Section~\ref{sec:algorithms}. Our contribution is
confined to software architecture, the GPU formulation of the factoring
stage, the verification methodology, and the measurements.

\section{Algorithms and their attribution}
\label{sec:algorithms}

We enumerate every algorithm the implementation relies upon, with its
source. Citations are attached in the code at the point of use, and we
preserve them when code is moved.

\subsection{Primality proving}

\begin{itemize}
  \item \textbf{The ECPP criterion.} Goldwasser and Kilian
    \cite{goldwasser1986}. If $E$ is an elliptic curve over
    $\mathbb{Z}/n\mathbb{Z}$ of order $m$, $q \mid m$ is prime with
    $q > (\sqrt[4]{n}+1)^2$, and $P \in E$ satisfies $[m/q]P \neq
    \mathcal{O}$ while $[q]([m/q]P) = \mathcal{O}$, then $n$ is prime.
  \item \textbf{The Atkin--Morain downrun.} Atkin and Morain
    \cite{atkin1993}; Cohen \cite[Algorithm 9.2.4, p.~474]{cohen1993}.
    Complex multiplication supplies curves of known order, and the
    criterion above reduces primality of $n$ to primality of the smaller
    $q$, recursively.
  \item \textbf{Miller--Rabin.} Miller \cite{miller1976} and Rabin
    \cite{rabin1980}, via GMP's \code{mpz\_probab\_prime\_p} with ten
    repetitions, used to reject composites before ECPP is attempted.
\end{itemize}

\subsection{Number-theoretic subroutines}

\begin{itemize}
  \item \textbf{Square roots modulo $p$.} Cohen
    \cite[Algorithm 1.5.1, p.~33]{cohen1993}, this being the
    Tonelli--Shanks method. Closed forms are used first for
    $p \equiv 3 \pmod 4$ and $p \equiv 5 \pmod 8$.
  \item \textbf{RESSOL.} Shanks' algorithm as described by Niven,
    Zuckerman and Montgomery \cite[pp.~110--114]{niven1991}; the
    implementation in this codebase is due to Scott Contini and serves as
    a final fallback.
  \item \textbf{Modified Cornacchia.} Cohen
    \cite[Algorithm 1.5.3, p.~36]{cohen1993}, after Cornacchia
    \cite{cornacchia1908}, solving $x^2 + |D|y^2 = 4p$.
  \item \textbf{Roots of a polynomial modulo $p$.} Cohen
    \cite[Algorithm 1.6.1, p.~37]{cohen1993}, an equal-degree splitting
    procedure in the manner of Cantor and Zassenhaus
    \cite{cantor1981}.
  \item \textbf{Lenstra's elliptic curve method.} Lenstra
    \cite{lenstra1987}; Cohen \cite[Algorithm 10.3.3, p.~488]{cohen1993}.
\end{itemize}

\subsection{Class polynomial construction}

\begin{itemize}
  \item \textbf{Hilbert class polynomials.} Cohen
    \cite[Algorithm 7.6.1, p.~415]{cohen1993}. Reduced forms $(a,b,c)$ of
    discriminant $D$ are enumerated, the $j$-invariant evaluated at each,
    and the linear and quadratic factors multiplied together. The
    $j$-invariant is built from the modular discriminant $\Delta$, itself
    obtained from a Dedekind eta partial sum.
  \item \textbf{Weber class polynomials.} Weber
    \cite[\S127]{weber1908}; Kaltofen and Yui \cite{kaltofen1991};
    Atkin and Morain \cite{atkin1993}. Weber polynomials have
    substantially smaller coefficients than their Hilbert counterparts
    and are correspondingly cheaper to construct.
  \item \textbf{Weber-to-$j$ conversion.} Kaltofen and Valente
    \cite{kaltofen1989}, whose Las Vegas primality test supplies the map
    from a Watson root to a Hilbert root, switching on $D \bmod 8$ and
    passing through unchanged when $D \equiv 0 \pmod{32}$.
  \item \textbf{Auxiliary predicates.} The functions denoted $G$, $I$,
    $J_1$, $K$, $L$, $M$ and $N$ implement helper predicates from the
    IEEE P1363 standard \cite{ieee1363}, following the presentations of
    Kaltofen--Yui and Atkin--Morain. Work on efficient curve generation
    over prime fields by Konstantinou, Stamatiou and Zaroliagis
    \cite{konstantinou2002} informs this portion of the implementation.
  \item \textbf{Transcendental constants.} $\pi$ is obtained at working
    precision by Hutton's Machin-like formula
    $\pi = 4(\arctan\frac{1}{2} + \arctan\frac{1}{3})$, evaluated with the
    library's own arctangent.
\end{itemize}

\subsection{GPU factoring}

\begin{itemize}
  \item \textbf{Montgomery curves and $XZ$ arithmetic.} Montgomery
    \cite{montgomery1987}, supplying the curve form
    $By^2 = x^3 + Ax^2 + x$, differential addition, and the ladder that
    together permit stage-one ECM without modular inversion in the inner
    loop.
  \item \textbf{Suyama's parametrisation.} Suyama, as reported by
    Montgomery \cite{montgomery1987} and Zimmermann and Dodson
    \cite{zimmermann2006}, generating curves with favourable torsion from
    a single parameter $\sigma$.
  \item \textbf{Montgomery modular multiplication.} Montgomery
    \cite{montgomery1985}, used throughout the kernel via CGBN's
    Montgomery primitives.
  \item \textbf{Stage two by prime continuation.} Montgomery's baby-step
    giant-step continuation \cite{montgomery1987}: primes $q = mD \pm j$ are
    detected by comparing $x$-coordinates of $[mD]Q$ and $[j]Q$, which needs no
    inversion in the loop.
  \item \textbf{FFT continuation.} Montgomery's polynomial stage two
    \cite{montgomery1992}, in the cuFFT backend: $\prod_{m,j} (X_m - x_j Z_m)$
    by polynomial products modulo $\prod_j (X - x_j)$ and multipoint
    evaluation by Bernstein's scaled remainder tree \cite{bernstein2004},
    with Newton reciprocals and Kronecker substitution as in von zur
    Gathen and Gerhard \cite{vzgathen2013}.
  \item \textbf{cuFFT.} NVIDIA's FFT library \cite{cufft}, for the
    double-precision convolutions behind every product above $16352$ bits.
  \item \textbf{CGBN.} The Cooperative Groups Big Numbers library of
    NVIDIA Labs \cite{cgbn}, which supplies fixed-width multiple-precision
    integer arithmetic distributed across cooperating threads.
  \item \textbf{GMP.} Granlund and collaborators \cite{gmp}, underlying
    all host arithmetic.
\end{itemize}

\section{From program to library}

\subsection{Three structural obstacles}

A program may print, may terminate, and may report its findings as text.
A library may do none of these. Each prohibition forced a specific change.

\paragraph{Output.} All fifty-one direct console writes were redirected
to a sink that accumulates output and delivers it to an optional callback
one complete line at a time. Line buffering is not incidental: the
original streams fragments of a single line across several functions, so
per-statement delivery would fragment the commentary. With no callback
installed the library is silent.

\paragraph{Termination.} The original called \code{exit(2)} from within
its modular polynomial arithmetic upon encountering a non-invertible
element. This is a proof of compositeness rather than an error, and it is
unacceptable for a library to terminate its host process on discovering
one. We compute $\gcd(y, m)$, which is precisely the nontrivial divisor
that the failed inversion exhibits, and throw an exception carrying it,
caught at the API boundary and reported as a composite result with the
factor attached.

\paragraph{Results as data.} The certificate is returned as an array of
structured steps. This additionally repairs an ambiguity in the original:
the command line tool returns exit status $2$ for a proved composite, for
exhaustion of discriminants, and for exceeding factoring bounds alike.
The library reports these separately while the tool retains its
historical behaviour.

\subsection{Configuration and concurrency}

The prover retains its configuration in file-scope globals. Rather than
thread a context through some 2\,600 lines of transcribed algorithms, we
serialise entry points with a mutex: concurrent \emph{calls} are safe but
not parallel. This is a limitation on inter-call concurrency only; within
a single call the prover does exploit several cores, as
Section~\ref{sec:parallel} describes. Callers wanting many proofs at once
should use separate processes, or the MPI driver.

One piece of state made that possible. Of everything the prover touches,
\code{gmp\_randstate\_t} is the only object it \emph{mutates} rather than
merely reads; the transcendental constants and configuration are
read-only once initialised. We verified the consequence rather than
assuming it: disjoint \code{mpz} and \code{mpf} operations across
thirty-two threads, while concurrently reading shared constants, produce
no discrepancy against a single-threaded reference. Giving each worker
its own random state is therefore sufficient, and it is installed in a
thread-local pointer that \code{rand2()} consults, because the call sites
are buried too deep in the elliptic curve and point-search code to have a
state threaded through them by hand.

The transcendental constants $\pi$, $e$ and $\ln 1.9$ are computed at
working precision, which at $10^4$ bits costs approximately $0.4$
seconds. We cache them by precision, so that a caller proving many
numbers pays this cost once. A separate entry point emits the banner and
computes these constants, allowing an interactive client to prompt for
input after the banner has appeared, as the original did.

\subsection{Interface}

The public interface comprises thirty-six exported symbols, all others
being hidden (thirty-two in the first release; the thread settings of
Section~\ref{sec:parallel} and the standalone GPU entry point of
Section~\ref{sec:cudaecm} added four). The header compiles as both C and C++, and a variant
accepting a GMP integer directly is declared when \code{<gmp.h>} precedes
it.

\section{The CGBN backend}

\subsection{Why only the factoring stage}

CGBN provides fixed-width integer arithmetic. Of the prover's four
layers, the class-polynomial layer is arbitrary-precision \emph{floating
point}, implemented as hand-rolled Taylor series with argument reduction;
the root-finding layer is sequential; and the downrun is inherently
serial, each step depending on its predecessor. Only elliptic curve
factorisation, in which every curve is an independent trial and a single
success suffices, admits a CGBN formulation.

\subsection{Kernel design}

Each CGBN instance carries one curve. Given $\sigma$, Suyama's
parametrisation gives $u = \sigma^2 - 5$, $v = 4\sigma$, initial point
$(u^3 : v^3)$, and
\[
  a_{24} = \frac{A+2}{4} = \frac{(v-u)^3(3u+v)}{16u^3v},
\]
requiring one modular inversion per curve. A failed inversion again
exhibits a factor, and is reported as such. Stage one then multiplies the
point by every prime power below the bound, using Montgomery's
differential addition and doubling in $XZ$ coordinates, so that no
inversion occurs in the inner loop. A nontrivial $\gcd$ of the final $Z$
coordinate with $n$ is a factor.

Instance width is fixed at compile time; the backend selects the
narrowest of $512$, $1024$, $2048$, $4096$, $8192$ and $16384$ bits that accommodates the input
(one $32$-bit limb is left spare, so the limits are $480$, $992$, $2016$, $4064$, $8160$ and
$16352$ bits). $16384$ is the practical ceiling: a CGBN instance spans at most one warp, and at
that width the hot kernels already use $236$--$248$ of the $255$ registers a thread may have.

\subsection{The kernel is latency bound}

Measured with stage one only (RTX 4080), the GPU's marginal cost at $B_1 = 10^4$
was $0.227$\,s for $256$ curves and $0.354$\,s for $4096$, and at $5\cdot10^4$
$1.119$\,s and $1.802$\,s: sixteen-fold increases in curve count cost $1.4\times$ in
wall time. Table~\ref{tab:kernel} repeats the measurement with the prover's current GPU
path, which adds stage two when stage one fails.

\begin{table}[t]
\centering
\begin{tabular}{rrrr}
\toprule
$B_1$ & 256 curves & 1024 curves & 4096 curves \\
\midrule
$10^2$ & 0.057\,s & 0.057\,s & 0.069\,s \\
$10^3$ & 0.056\,s & 0.057\,s & 0.069\,s \\
$10^4$ & 0.483\,s & 0.489\,s & 0.337\,s \\
$5\times10^4$ & 1.233\,s & 1.209\,s & 1.504\,s \\
\bottomrule
\end{tabular}
\caption{GPU marginal cost (a CUDA run minus a host-only run of the same work; 46-digit
  modulus with a 20-digit factor; RTX PRO 4500, stage one and two). Cost tracks $B_1$;
  at $B_1 = 10^4$ more curves were \emph{cheaper}, because stage one then more often found
  the factor and stage two was skipped.}
\label{tab:kernel}
\end{table}

Cost tracks $B_1$, which determines the length of the serial chain each
curve must traverse, and is nearly independent of curve count until the
device saturates, at about one wave of curves (Section~\ref{sec:cudaecm}). Curves are therefore almost free. This is the shape ECM
rewards, and it is the reason the backend defaults to thousands of
concurrent curves against the host implementation's twenty per bound.

\subsection{Dispatch policy: a negative result}

Our initial implementation dispatched to the GPU whenever the CUDA
backend was selected. This made whole proofs approximately twice as slow.
The cause is that most factoring calls in an ECPP downrun exist to reject
candidate curve orders, and are \emph{expected} to fail; a fixed
per-launch cost paid on every such rejection is pure loss. The backend
consequently attempts the host implementation first and escalates to the
GPU only when the host has failed \emph{and} the cofactor exceeds a
configurable width, $192$ bits by default. Under this policy enabling the
backend is neutral where it cannot help and substantially profitable
where it can.

\subsection{A standalone, resumable GPU ECM}
\label{sec:cudaecm}

The dispatch policy above suits a prover, whose factoring calls are meant to fail fast. A program
that \emph{factors} wants the opposite: the GPU as an engine of its own, run beside GMP-ECM on the
host. The first such client, the trial-division factorizers of the \code{primefactors} project,
gained almost nothing from \code{gmpecpp\_ecm\_factor()}. Profiling showed why: each call
spent $1.36$\,s of $2.04$\,s at $B_1 = 10^4$, and $6.99$\,s of $8.92$\,s at $5\cdot10^4$, in the
host attempt rather than the kernel ($0.34$\,s and $1.52$\,s); $4096$ curves filled the device only
part way; and stage one alone, with $B_1 \le 10^5$, made each GPU curve 6--20$\times$ less likely
to succeed than a GMP-ECM curve with stage two.

\code{gmpecpp\_cuda\_ecm()} --- taking $n$, $B_1$, $B_2$, a curve count, a seed and a stop
callback --- addresses each point:
\begin{itemize}
  \item \textbf{GPU only.} No host attempt and no global lock: concurrent calls, on one device or
    several, are safe. The call runs on the caller's current device.
  \item \textbf{One wave of curves.} \code{curves = 0} asks the occupancy API how many curves the
    device runs at once ($11\,808$ at $512$ bits, $5\,904$ at $1024$, $2\,952$ at $2048$, $1\,968$ at $4096$, $984$ at $8192$, $656$ at $16384$ on an
    RTX PRO 4500 Blackwell). Half a wave finds $21\%$ less per second; two waves find the same at
    twice the time per call.
  \item \textbf{Resumable launches.} $X$, $Z$ and $a_{24}$ of every curve stay in device memory,
    and stage one runs as launches of about $0.25$\,s, the slice of prime powers adapted from the
    previous launch's time. A \code{stop} callback is polled between launches, and $B_1$ is no
    longer bounded by how long one launch may take.
\end{itemize}
A call now costs $0.65$\,s at $B_1 = 10^4$ and $2.6$\,s at $5\cdot10^4$.

\subsection{Stage two on the GPU}

Let $Q$ be a curve's point after stage one and $D = 210$. Every prime $q \in (B_1, B_2]$ is
$mD \pm j$ for one $m$ and one of the $24$ odd $j < D/2$ coprime to $D$, and $[q]Q = O \pmod p$
exactly when $x([mD]Q) \equiv x([j]Q) \pmod p$, since $x$ does not see the sign of $j$. Each
curve stores its $24$ baby points $[j]Q$, computed by differential additions from $[1]Q$ and
$[2]Q$ and \emph{normalized} once to $x_j = X_j/Z_j$ (one inversion each; a failed inversion
exposes $\gcd(Z_j, n)$ as a factor). The giant points $[mD]Q$ advance by one differential
addition each, $[(m+1)D] = [mD] + [D]$ with difference $[(m-1)D]$, and for every pair $(m, j)$
that covers a prime the accumulator is multiplied by $X_m - x_j Z_m$; $\gcd(\mathit{acc}, n)$
ends the stage. That is two modular multiplications per prime, against about sixteen per unit
of $B_1$ in stage one. Which pairs to take does not depend on the curve, so the host builds them
once, as one $24$-bit mask per $m$ from a segmented sieve, and every curve reads the same mask.

Before the normalization each pair cost three multiplications ($X_m Z_j - X_j Z_m$); the
normalized form left every hit unchanged on identical curves and made stage two $1.4\times$
faster. Correctness was checked statistically, since a single run cannot tell a working stage
two from a lucky stage one: on a 20-digit factor at $B_1 = 10^4$ the per-curve success rate rose
from $0.00068$ to $0.0099$ with $B_2 = 50\,B_1$, against $0.0129$ for GMP-ECM with the same $B_2$.
An environment variable, \code{GMPECPP\_CUDA\_ECM\_STATS}, reports the number of successful curves
per stage, which is how these rates were measured.

\begin{table}[t]
\centering
\begin{tabular}{lrrrr}
\toprule
 & $25\,B_1$ & $50\,B_1$ & $100\,B_1$ & $200\,B_1$ \\
\midrule
20 digits, $B_1 = 10^4$ (three mults/prime) & 108 & 120 & \textbf{132} & 109 \\
30 digits, $B_1 = 2.5\cdot10^5$ (two mults/prime) & --- & \textbf{0.47} & 0.41 & 0.42 \\
\bottomrule
\end{tabular}
\caption{Factors found per second, one wave of curves, by the stage-two bound $B_2$. Stage two
  alone raised the 20-digit rate from 12 to 132.}
\label{tab:stage2}
\end{table}

Table~\ref{tab:stage2} chooses $B_2$: $100\,B_1$ at $B_1 = 10^4$ and $50\,B_1$ from $5\cdot10^4$ up.
The best $B_1$ for each factor size is GMP-ECM's own ($5\cdot10^4$ for 25 digits, $2.5\cdot10^5$ for
30 at $1.7$\,s per find, $10^6$ for 35 at $17$\,s per find), and there one GPU finds about
$9\times$ as many 30-digit factors per second as GMP-ECM on thirty-one host threads with its default
$B_2$, and 4--9$\times$ as many 35-digit ones. Run concurrently with GMP-ECM in the
\code{primefactors} factorizers, eight 400-bit inputs with 18--35-digit factors took $171$\,s
against $202$--$329$\,s on the host alone, the GPU winning most of the hardest.

The prover itself now uses the same kernels when it escalates to the GPU: stage one and stage
two, $B_2$ by the rule above, $B_1$ still clamped to $[10^3, 10^5]$ and $4096$ curves by default,
the host attempt still first. Over the thirty-six CUDA proofs of the sweep in
Section~\ref{sec:results} (90 to 180 digits, two seeds), the total fell from $145.9$\,s to
$117.9$\,s, every proof completed and every certificate verified on both backends. The gain is
where the GPU is reached --- 120 digits $4.8 \to 2.9$\,s, 150 digits $9.9 \to 3.2$\,s --- and
small inputs are unchanged. A full wave of curves gave $121.5$\,s, no better than $4096$.

\section{The cuFFT backend: ECM beyond the CGBN widths}
\label{sec:fft}

CGBN stops at $16384$-bit instances ($n$ up to $16352$ bits), and schoolbook multiplication
inside one warp would be the wrong algorithm well before any larger width could be compiled. A
second library, \code{libgmpecmfft}, with its own command line tool \code{gmpecmfft}, factors
numbers of any size. Its entry point \code{gmpecmfft\_ecm()} takes the same arguments as
\code{gmpecpp\_cuda\_ecm()} and simply calls it up to $16352$ bits; above that, every modular
multiplication is a floating-point convolution computed with NVIDIA's cuFFT \cite{cufft}. Both
paths draw $\sigma = 6 + \mathrm{random}(2^{30})$ from the seed in the same way, so a given seed
builds the same Suyama curves on either path, and the same curves as GMP-ECM's
\code{-sigma 0:}$\sigma$. That correspondence is the basis of every test below. The library links
\code{libcufft} rather than \code{libcufftw}: the FFTW-compatible interface takes host arrays and
would copy every operand across the bus on every call.

\subsection{Modular arithmetic by convolution}

A residue is a signed integer held as $L$ \emph{balanced} digits, each in
$[-2^{b-1}, 2^{b-1})$, in 32-bit integers, and a \emph{slot} holds one
residue for each of $C$ curves. All curves walk the same scalar, so every operation is on whole
slots and the curves move in lockstep, exactly as the CGBN threads do. A product of two slots is
a batched real-to-complex transform of length $N = 2L$, a pointwise product and an inverse
transform; each coefficient is then rounded to an integer and carries are propagated. Balanced
digits bound a convolution coefficient by $L\,2^{2b-2}$, a quarter of what unsigned digits
allow, so the digit size is the largest $b$ with $2b + \log_2 L \le 46$, which keeps every exact
coefficient below $2^{44}$, well inside the $53$-bit double mantissa: one bit per digit more than
unsigned digits, and $L$ correspondingly shorter. The rounding error is not assumed small but
measured on every product; the worst seen across the test sizes is about $0.0005$ (lower than
with unsigned digits, despite their larger size, since balanced sums typically stay far below
their bound), and a run whose error exceeds $0.25$ stops with a precision error rather than
continuing on arithmetic it cannot trust.

$L$ must be even and $7$-smooth (for the reason below), but is not simply rounded up to the
next such number: cuFFT's speed depends on a length's factorization more than on its size. At
$20\,000$ bits, stage one ran $15\%$ faster with $L = 1280 = 2^8 \cdot 5$ than with the
shorter $L = 1200 = 2^4 \cdot 3 \cdot 5^2$, and the best length differed with the number of
curves. So the library times the candidate lengths up to a quarter beyond the minimum on the
device itself when it sets up, and keeps the fastest; a larger $L$ only makes $R$ larger.

Reduction is Montgomery's \cite{montgomery1985}, with $R = 2^{bL} > 32n$. Bounds are on
absolute values: the reduction's multiplier $m$ is balanced, $|m| \le R/2$, so inputs below $4n$
give outputs below $n$ without any final subtraction, and an addition is $a + b$ and a
subtraction $a - b$, neither needing a comparison or an offset. Residues are never reduced to
$[0, n)$; ECM only ever asks for $\gcd(Z, n)$, which does not care. Carry
propagation runs one thread block per number, each thread rippling a chunk and chunk carries
passing upward until none remain.

Of REDC's two products, $m = (T \bmod R)\,n' \bmod R$ and $mn$, only the high half of the
second is used, and it can be had at half length, the device GMP's own \code{mpn\_redc\_n}
uses. Write $T = T_{\mathrm{hi}}R + T_{\mathrm{lo}}$. Since $mn \equiv -T_{\mathrm{lo}}
\pmod R$, and $R \equiv 1 \pmod{R - 1}$,
\[
\frac{T + mn}{R} \equiv T_{\mathrm{lo}} + T_{\mathrm{hi}} + (mn \bmod (R - 1))
\pmod{R - 1},
\]
the carry between the halves cancelling; and as the result is below $2n < R - 1$, that residue
is the result itself. $mn \bmod (R - 1)$ is a cyclic convolution of length $L$ rather than a
linear one of length $2L$, and the sum is normalized with the carry out of the top digit wrapped
around to the bottom. Every multiplication pays one REDC, so this shortens stage one and stage
two alike. The full-length reduction's exactness check, a low half that must come out zero, has
no counterpart here; instead the digits above $2n$'s length must be zero, which a wrong $m$
violates.

Operands are either slots or \emph{stored transforms}. Within one batched call, each distinct
slot is transformed once, however many products use it; values that never change --- the base
point and $a_{24}$ in stage one, the baby points in stage two --- are transformed once for the
whole run. A ladder step is three rounds of $4$, $4$ and $3$ multiplications, and its first
round needs only four forward transforms for six operands. Each step, for either bit value, is
captured once as a CUDA graph and replayed, which removes launch overhead from a loop of tens of
thousands of steps. Stage one is one Montgomery ladder over $k = \prod_{p^e \le B_1} p^e$, which
reaches the same point as CGBN's ladder per prime power.

\subsection{Two stage twos}

The \emph{per-prime} stage two is the CGBN algorithm of the previous section transcribed: $D =
210$, the host's prime masks, baby points normalized to $x_j$ (a batch inversion: the products
of the $Z_j$ up a product tree on the GPU, one GMP inversion per curve at the root, and back
down, a child's inverse being its parent's times its sibling's product), and
$X_m - x_j Z_m$ per prime. It visits exactly the pairs CGBN visits, so hit counts compare
exactly.

The \emph{polynomial} stage two is Montgomery's FFT continuation \cite{montgomery1992}, the
method GMP-ECM uses. With $D = 2310$ and the $k = \varphi(D)/2 = 240$ babies $x_j$, form
\[
F(X) = \prod_j (X - x_j), \qquad
G_i(X) = \prod_{m \in \text{block } i} (Z_m X - X_m),
\]
the giants taken $k$ at a time. Accumulating $H \leftarrow H\,G_i \bmod F$ over the blocks and
finally evaluating $\prod_j H(x_j)$ yields the product of $X_m - x_j Z_m$ over \emph{every} pair,
that is over every number coprime to $D$ in range, a superset of the primes. The work per block
is a product tree for $G_i$ and one multiplication and remainder modulo $F$, so the cost per
prime falls as $k$ grows. The final evaluation is Bernstein's \emph{scaled} remainder tree
\cite{bernstein2004} over $F$'s product tree: a node $P$ carries the first $\deg P$
coefficients of $(H \bmod P)/P$ as a series in $1/X$, a child $L$ of $P = LR$ obtains its own
as a middle product of its parent's with $R$, and at a leaf $X - x_j$ the single coefficient is
$H(x_j)$. Only the root needs a reciprocal, $1/\mathrm{rev}(F)$, which the blocks already use
(Newton iteration \cite{vzgathen2013}); a plain remainder tree would need one per node and
about four times the reductions. Larger choices ($D =
30030$, $k = 2880$) and a smaller one ($D = 210$, $k = 24$) are available. The polynomial form is
used when $D \le B_1/2$ --- keeping the first giant past $[D]Q$, since a differential addition
with difference $O$ is degenerate --- and $B_2 - B_1 \ge Dk/2$; otherwise the per-prime form
runs. The last block is filled rather than truncated, which costs nothing and makes the
effective $B_2$ slightly larger than requested.

Polynomial products use Kronecker substitution: each curve's coefficients are packed $2L$ digits
apart into one long integer, multiplied by a single large transform, and every coefficient of
the product is read back as a $2L$-digit value and reduced by the same REDC as a scalar product.
A coefficient of a product of degree-$k$ polynomials sums up to $k + 1$ terms, so this stage
runs on its own instance of the arithmetic with smaller digits ($2b + \log_2 L + \lceil\log_2(k+1)\rceil \le
46$) and a larger $R > 32(k+1)n$; stage one's instance is freed first, since each sizes its
buffers from free device memory, and the stage-one point moves between the two Montgomery forms
through the host. Remainders $A - FQ$ stay below $3n$ and are only ever used as product inputs,
so coefficient bounds do not grow down a tree. In a block's
$H G_i \bmod F$ the product $HG_i$ is instead kept exact: only its top $k$ coefficients are
reduced, to form the quotient $Q$, and each low coefficient is reduced once as the difference
of the exact $HG_i$ and $FQ$, which $R > 32(k+1)n$ accommodates --- three reductions per
coefficient of $F$'s degree where
the general remainder takes four. Reversal and truncation of polynomials only reorder slot
pointers.

Memory per curve grows with $k$: at $D = 30030$ and $20\,000$ bits a curve needs about a
gigabyte, and a $32$\,GB device holds sixteen. The stage therefore runs on \emph{groups} of curves,
one after another, each with its own instance of the arithmetic: the largest $D$ that fits for
some group size, the group as large as fits. Since even one curve's Kronecker transforms are
large, the cost per curve hardly depends on the group size, and a large $D$ in small groups
beats a small $D$ in one.

\subsection{Validation}

Nothing in this backend is checked by comparing timings. \code{make check-fft} establishes,
in a few minutes:
\begin{itemize}
  \item \textbf{Arithmetic against GMP.} Rounds of products, squares, sums and differences on
    random residues mod $n$, and, up to $10^5$ bits, polynomial products, product and remainder
    trees, Newton inverses and remainders: no wrong result from $200$ to $10^6$ bits. A
    deliberately wrong expectation is reported as $528$ errors, so the check can fail.
  \item \textbf{Against CGBN.} On a 2000-bit input, $256$ curves, $B_1 = 10^4$, $B_2 = 10^5$
    (the per-prime stage two at that $B_2$), the FFT path finds a factor on exactly the same
    number of curves as CGBN in stage one and in stage two.
  \item \textbf{The polynomial stage two against a host replay.} With
    \code{GMPECMFFT\_S2\_CHECK=1} two curves' polynomial stage two is recomputed with GMP ---
    the same ladders, chains and formulas, hence the same projective coordinates --- and the
    product over all pairs must agree exactly, for $D = 2310$ and for $D = 210$ forced.
  \item \textbf{Curve groups.} Stage two run in groups of seven curves finds a factor on
    exactly the same curves as in one group of $128$.
  \item \textbf{Against GMP-ECM.} On a 20\,000-bit input, $256$ curves with the same
    $\sigma$ and $B_1 = 5000$, stage one succeeds on the same number of curves under both
    programs.
\end{itemize}

No committed harness measures this backend's speed, so no figures are given here. Profiles on
the RTX PRO 4500 show cuFFT's double-precision kernels dominating; that card's FP64 rate is the
bound, and the backend is meant for devices with a full-rate FP64 unit.

\section{Precision as the dominant parameter}

Table~\ref{tab:phases} reports per-phase exclusive time at two working
precisions.

\begin{table}[t]
\centering
\begin{tabular}{lrrr}
\toprule
Phase & 30 digits & 60 digits & 100 digits \\
\midrule
\multicolumn{4}{l}{\emph{At $10^3$ bits}} \\
Elliptic curve factorisation & 51.0\% & 98.2\% & 99.5\% \\
Hilbert class polynomials    & 35.6\% & 1.5\%  & 0.1\% \\
Total                        & 0.062\,s & 3.859\,s & 19.42\,s \\
\midrule
\multicolumn{4}{l}{\emph{At $10^4$ bits}} \\
Elliptic curve factorisation & 1.0\%  & 28.9\% & 89.8\% \\
Hilbert class polynomials    & 84.3\% & 67.5\% & 7.8\% \\
Total                        & 3.015\,s & 12.530\,s & 21.09\,s \\
\bottomrule
\end{tabular}
\caption{Phase distribution against working precision.}
\label{tab:phases}
\end{table}

The factoring rows read across as $0.032$, $3.790$, $19.320$ seconds at
$10^3$ bits against $0.032$, $3.619$, $18.934$ seconds at $10^4$: the
absolute cost of factorisation is unchanged, since precision does not
enter integer arithmetic. What changes is the cost of everything around
it.

We therefore reduced the default precision to $10^3$ bits, after
verifying across the thirty-three primes of the committed benchmark
set, from $20$ to $180$ decimal digits, that
every input still proves, that no loss-of-precision event occurs, and
that no proof becomes slower. Observed speedups range from $48.5\times$ at
thirty digits to $1.1\times$ beyond one hundred, with median
$1.70\times$. The failure mode this
trades against is bounded: insufficient precision causes a
consistency check on the constant coefficient to fail, whereupon the
discriminant is abandoned and the sweep continues. It costs time, not
correctness.

\section{Parallelism}
\label{sec:parallel}

\subsection{Where it exists, and where it does not}

An ECPP downrun is sequential by construction: step $i+1$ establishes the
primality of the prime step $i$ produced, and no amount of hardware
shortens that chain. What the algorithm does offer is the finer axes of
Table~\ref{tab:axes}, which do not all suit the same tool.

\begin{table}[t]
\centering
\begin{tabular}{llrl}
\toprule
Site & Independent? & Runtime & Mechanism \\
\midrule
ECM curves & fully; first success ends it & 52--99.5\% & POSIX threads \\
Discriminant sweep & until the first success & $\approx$100\% & POSIX threads \\
Whole numbers & fully & --- & MPI \\
$m$ candidates per $D$ & fully & within ECM & unused \\
$j(\tau)$ per reduced form & fully, then a reduction & 0.1--85\% & unused \\
Root finding & two subtrees & $<$1\% & not worthwhile \\
The downrun & \textbf{no} & --- & --- \\
\bottomrule
\end{tabular}
\caption{Parallelism available in the prover.}
\label{tab:axes}
\end{table}

\subsection{Elliptic curve trials}

Curves are independent and a single success ends the search, so the
twenty attempts per bound become twenty concurrent attempts. Each worker
owns a random state; nothing else the inner loop touches is mutable
shared state.

\begin{table}[t]
\centering
\begin{tabular}{rrr}
\toprule
Threads & Total & Speedup \\
\midrule
1  & 171.6\,s & 1.00$\times$ \\
2  & 135.0\,s & 1.27$\times$ \\
4  &  54.9\,s & 3.12$\times$ \\
8  &  41.9\,s & 4.09$\times$ \\
16 &  31.9\,s & 5.38$\times$ \\
32 &  28.7\,s & \textbf{5.98$\times$} \\
\bottomrule
\end{tabular}
\caption{ECM curve threading over the eighteen primes of ninety digits or
  more in the benchmark set, thirty-two cores.}
\label{tab:threads}
\end{table}

The aggregate in Table~\ref{tab:threads} is sublinear and flattens beyond
sixteen threads, yet individual inputs routinely exceed the thread count:
one ninety-digit prime falls from $11.63$ to $2.52$ seconds on two
threads, and a hundred-and-forty-digit one from $52.63$ to $2.72$ seconds
on four. Both follow from ECM being a Las Vegas algorithm. Concurrent
curves finish in the time of the luckiest rather than the average, giving
large gains exactly where the sequential search was unlucky and none
where it was not; averaged over a sample that is the column above. The
flattening arrives once trials past roughly twenty curves per bound are
purely speculative.

A methodological caution belongs here. An earlier measurement over five
primes reported $6.96\times$. That figure was not wrong for its sample,
but a small sample of a high-variance quantity flatters whichever tail it
catches, and $5.98\times$ over eighteen primes is the defensible number.

\subsection{The discriminant sweep}

Each discriminant is an independent trial and the work beneath one is
essentially a whole proof step, so a speculative parallel sweep covers
almost the entire computation. Two design points make it usable rather
than merely faster.

\paragraph{Determinism.} Every attempt seeds its own random stream from
the discriminant it tests, making the attempt a pure function of
$(D, N, \text{seed})$. Results are therefore independent of scheduling:
thirty-two threads reproduce two threads byte for byte. This is unusual
for speculative parallelism and worth its small cost.

\paragraph{The smallest discriminant still wins.} Workers draw
monotonically increasing indices; a success at index $k$ halts new work
beyond $k$ while in-flight attempts below it finish, and the smallest
successful index is taken. The sweep therefore selects the discriminant
the sequential walk would have reached. This matters because a larger
$|D|$ means a larger class number $h(D)$ and a costlier proof step:
accepting whichever success arrived first would often choose a worse
curve and surrender the speedup.

\begin{table}[t]
\centering
\begin{tabular}{rrr}
\toprule
Threads & Total & Speedup \\
\midrule
1  & 172.3\,s & 1.00$\times$ \\
2  & 133.6\,s & 1.29$\times$ \\
4  & 122.3\,s & 1.41$\times$ \\
8  & 119.9\,s & 1.44$\times$ \\
16 & 118.7\,s & 1.45$\times$ \\
32 & 118.5\,s & \textbf{1.45$\times$} \\
\bottomrule
\end{tabular}
\caption{Discriminant sweep threading, same inputs as
  Table~\ref{tab:threads}.}
\label{tab:sweep}
\end{table}

At $1.45\times$ against $5.98\times$ on identical inputs this is
substantially weaker, for structural reasons rather than defects of
implementation. Most inputs succeed on one of the first few
discriminants, leaving little sweep to parallelise; and because the
smallest success is retained, every other worker's output is discarded.
That discarding is the design, not waste to be optimised away: it is what
buys the cheaper curve.

The distribution is accordingly bimodal. Where the sweep is long the gain
is large -- $52.98$ to $23.68$ seconds at a hundred and forty digits,
$19.08$ to $10.14$ at a hundred -- while inputs whose first discriminant
already worked lose a few to twenty percent to thread creation. An
earlier version lost far more, turning $0.66$ seconds into $2.29$; the
smallest few discriminants are now probed on the calling thread before
any fan-out, which is free in the parallel case.

\subsection{Distributed batches, and why neither tool suffices}

Whole numbers are embarrassingly parallel with no communication until the
end, which is what MPI is for. Work is dispatched dynamically rather than
partitioned in advance, because proof times vary by two orders of
magnitude between neighbouring inputs and a static division strands most
ranks behind one straggler.

\begin{table}[t]
\centering
\begin{tabular}{lrr}
\toprule
Configuration & Wall & Speedup \\
\midrule
1 rank, 1 thread   & 53.78\,s & 1.00$\times$ \\
5 ranks, 1 thread  & 20.47\,s & 2.63$\times$ \\
9 ranks, 1 thread  & 20.45\,s & 2.63$\times$ \\
5 ranks, 4 threads &  5.68\,s & 9.47$\times$ \\
9 ranks, 4 threads &  4.89\,s & \textbf{11.00$\times$} \\
\bottomrule
\end{tabular}
\caption{Ten primes of ninety to a hundred and ten digits.}
\label{tab:mpi}
\end{table}

Table~\ref{tab:mpi} carries the argument for combining the mechanisms
rather than choosing between them. MPI plateaus at $2.63\times$ and
ranks beyond five change nothing: a batch's wall time is bounded below by
its slowest single proof, and no number of ranks divides one proof.
Threads do, because they operate inside it, so the hybrid reaches
$11.0\times$ where either tool alone stalls.

\subsection{One number's curves across GPUs}
\label{sec:fftmpi}

The factorizer of Section~\ref{sec:fft} offers a different axis to distribute. Its ECM curves
are independent trials, and the first to find a factor ends the search, so a cluster's GPUs
can share one number's curves while almost nothing crosses the network: the number, the
bounds, and which curves to run. That makes curves, not the arithmetic of a single curve, the
unit to distribute. Splitting one curve's transforms across nodes would make every
multiplication wait on the network, and gain nothing in capacity, since even at $10^6$ bits a
curve's values occupy a few hundred kilobytes each. For the same reason there is no bulk data
for a GPU collective library such as NCCL to move, and the tool is written against MPI alone.

\code{gmpecmfft-mpi} gives rank $0$ to a coordinator that uses no GPU; every other rank is a
worker on one GPU, chosen by its index among the workers on its node. A batch is one call of
\code{gmpecmfft\_ecm()} on a worker's GPU, by default one wave of curves. As in the prover's
batch mode, work is handed out on request rather than divided in advance: a worker's result
is also its request for the next batch, so no GPU waits on another. Batch $i$ draws its curves
from seed $S + i$. The batches are therefore distinct, and any one of them, in particular the
one that found a factor, is reproduced exactly by a single-GPU run with that seed.

Stopping has to be prompt, since a batch at large bounds runs for minutes. When a factor
turns up, after a requested number of batches, or on an interrupt, the coordinator sends every
worker a single stop message. A busy worker notices it through the library's stop callback,
which the library polls about four times a second and which here probes for that message
without consuming it. The worker abandons its batch, reports once more, and then receives the
stop. Each worker thus gets exactly one stop and sends exactly one final report, which leaves
no message unmatched when the job ends. A worker that cannot run, or whose batch fails, is
retired and its batch given to another; the run fails only if every worker does.

The expected time to a factor falls in proportion to the number of GPUs only while the cluster
needs more than one batch per GPU to run the expected number of curves for that factor's size.
Beyond that point more GPUs no longer help, since a batch cannot finish faster than one curve
does. The GPUs should also be of one kind. Batches are of equal size, so with a mix, the
faster GPUs finish their share and then sit idle at the end while the slower ones complete
theirs.

The tool has been checked with two workers sharing one GPU, not yet across nodes. The factor
came from exactly the batch that single-GPU runs with the same seeds predict, both on the CGBN
path at $2\,000$ bits (a test \code{make check-fft} now runs) and on the FFT path at
$20\,000$ bits. An interrupt delivered to the coordinator alone stopped workers that were never
signalled in the middle of their batches. No multi-GPU timings are reported, since none has been
measured.

The same batches can also be run from a Slurm job array, by \code{tools/gmpecmfft-slurm}, with
no MPI and no coordinator. Each array task holds one GPU, and task $t$ of $T$ runs batches $t$,
$t + T$, $t + 2T$, \ldots, each batch exactly as the MPI tool would run it. The only shared state
is a run directory on a filesystem every node sees. The first task to find a factor claims it by
creating a directory there, which is atomic, records the factor, leaves a stop file, and cancels
the rest of the array. Running tasks then stop on the scheduler's termination signal; where
cancelling is not permitted, they stop within a quarter second of the stop file appearing,
which the single-GPU tool now polls on request. This needs nothing but the scheduler and a
shared filesystem, and tasks can start, finish and be requeued independently. The price is that
the division of work is fixed in advance: a task on a slower GPU still has its full share of
batches to run, so the requirement of one kind of GPU is stricter here than with the
coordinator's on-demand dispatch. The scheduler itself was not available on the development
machine, so the array tasks were checked by running two of them by hand: the factor came from the
predicted batch, and the other task stopped in the middle of its batch through the stop file.

\subsection{Two defects, and what found them}

Both are recorded because each was invisible to the obvious test.

\paragraph{Flags surviving a step.} The success flags live in the downrun
driver but were cleared inside the loop body that became the
per-discriminant function. A sweep finding nothing therefore inherited
the previous step's success and continued with its stale curve; a
hundred-and-forty-digit proof ran two hundred and fifty-seven steps and
emitted a certificate failing every check. The regression matrix of four
small inputs could not detect this, since the fault requires a downrun
long enough for some sweep to fail. That matrix now has a mode replaying
every prime in the benchmark set, which does.

\paragraph{\code{volatile} mistaken for synchronisation.} The shared
found-flags were declared \code{volatile}, read every iteration and
written once. \code{volatile} orders nothing and the pair is a data race
whatever a given architecture does in practice; ThreadSanitizer
identified it. Both are atomics now, and the sanitiser reports no races
across either threaded stage.

\paragraph{Defaults.} Both mechanisms default to a single thread and
reproduce the historical program byte for byte. Above one they alter
which factor is discovered, so certificates differ legitimately and are
checked with the verifier of Section~\ref{sec:verify} rather than
compared.

\section{GMP or CGBN: a measured comparison}
\label{sec:gmpcgbn}

A natural question, once part of the arithmetic runs on the GPU, is whether the rest should
follow: could CGBN take over every place where the host uses GMP, in the prover and in the
factorizer? This section answers it with measurements. All of them were taken on the
development machine, an Intel Core i9-13900K (8 performance cores with two threads each and 16
efficiency cores, 32 hardware threads) with the RTX PRO 4500, CUDA 13.2, GMP 6.3.0 and Open
MPI 5.0.5, and every table comes from a committed harness (\code{make bench-cgbn-latency},
\code{bench-gmp-profile}, \code{bench-mpi-cgbn}).

\subsection{The cost of one operation}

Table~\ref{tab:cgbnlat} compares one GMP operation with CGBN in two ways: as a drop-in
replacement for that single call (copy the operands in, launch one instance, synchronize, copy
the result out), and as one of 4096 independent instances in a single launch. Every CGBN result
was checked against GMP.

\begin{table}[t]
\centering
\begin{tabular}{lrrrr}
\toprule
Operation & GMP & CGBN, one & CGBN, 4096 at once & \\
\midrule
$a^e \bmod n$, 512 bits  & 49\,$\mu$s   & 1439\,$\mu$s ($29\times$ slower) & 0.83\,$\mu$s & ($59\times$ faster) \\
$a^e \bmod n$, 1024 bits & 321\,$\mu$s  & 4774\,$\mu$s ($15\times$ slower) & 3.14\,$\mu$s & ($102\times$ faster) \\
$a^e \bmod n$, 2048 bits & 2157\,$\mu$s & 9210\,$\mu$s ($4.3\times$ slower) & 17.4\,$\mu$s & ($124\times$ faster) \\
$ab \bmod n$, 512 bits   & 0.15\,$\mu$s & 26.2\,$\mu$s ($172\times$ slower) & 0.013\,$\mu$s & ($12\times$ faster) \\
$ab \bmod n$, 1024 bits  & 0.43\,$\mu$s & 26.2\,$\mu$s ($61\times$ slower) & 0.019\,$\mu$s & ($22\times$ faster) \\
$ab \bmod n$, 2048 bits  & 1.42\,$\mu$s & 26.3\,$\mu$s ($19\times$ slower) & 0.036\,$\mu$s & ($40\times$ faster) \\
\bottomrule
\end{tabular}
\caption{One GMP operation (one CPU thread) against one CGBN instance, round trip included, and
  against 4096 CGBN instances in one launch (per operation). Exponents are full size; CGBN uses
  8, 16 and 32 threads per instance at 512, 1024 and 2048 bits.}
\label{tab:cgbnlat}
\end{table}

The two columns tell opposite stories. A single CGBN instance occupies between 8 and 32
threads of one streaming multiprocessor, so it has none of the GPU's parallelism to draw on, and
every call pays a fixed round trip of about $26\,\mu$s, the floor visible in the
multiplication rows. Replacing a GMP call one for one is therefore $4$ to $172$ times slower.
Thousands of independent operations of one width are another matter: they fill the device, and
CGBN is then $12$ to $124$ times faster than GMP. CGBN only pays where the work comes in large
batches.

\subsection{Where the host's time goes}

Table~\ref{tab:gmpshare} bounds what any move of GMP work to the GPU could save: the share of
host time that falls in \code{libgmp}. A thread blocked in a CUDA synchronization spins in the
driver, so the \code{libcuda} share is time spent waiting for the GPU.

\begin{table}[t]
\centering
\begin{tabular}{lrrrrr}
\toprule
Run & Wall & \code{libgmp} & gcd, division & \code{mpf} & \code{libcuda} \\
\midrule
Prover, 140 digits, host only & 67.6\,s & 92\% & 76\% & 0.0\% & --- \\
Prover, 140 digits, \code{-C} & 2.9\,s  & 65\% & 52\% & 0.1\% & 24\% \\
\code{gmpecmfft}, 2\,000 bits (CGBN path) & & 0.04\% & & & 90\% \\
\code{gmpecmfft}, 20\,000 bits (cuFFT path) & & 0.06\% & & & 89\% \\
\bottomrule
\end{tabular}
\caption{Share of host time by library, and for the prover the part of \code{libgmp} spent in
  gcd and division routines and in \code{mpf} (\code{perf}, \code{cpu-clock} at 1999\,Hz). The
  prover runs one thread, seed $1$; \code{gmpecmfft} runs $2952$ and $64$ curves,
  $B_1 = 10^4$, $B_2 = 10^6$.}
\label{tab:gmpshare}
\end{table}

The factorizer has almost nothing to give: its host GMP work, curve setup and one gcd per curve,
is at most $0.06\%$ of its time, and above $16352$ bits CGBN could not even hold the numbers.
In the prover, the host-only run is dominated by gcd and division routines
(\code{mpn\_hgcd2}, \code{div2}): the prover's point arithmetic is affine and pays a modular
inversion for every point addition, both in its elliptic curve method and in the point checks of
each downrun step. The elliptic curve method's share of that is exactly what the CGBN backend
replaces, with inversion-free Montgomery curves run thousands at a time; it is why \code{-C}
takes the proof from $67.6$ to $2.9$\,s. With \code{-C} the remaining host time is still
mostly the same affine arithmetic (gcd and division, $52\%$; the class polynomials' \code{mpf}
arithmetic is only $0.1\%$ at the default precision). It belongs to work that is small and
sequential: the host's own first attempt at every factorization, which the dispatch policy keeps
on the host because most attempts are meant to fail quickly (escalating all of them was measured
to make proofs slower), and the downrun's point checks, a short chain per step. Neither offers
the thousands of independent operations Table~\ref{tab:cgbnlat} requires. The GMP work that
batches has already moved; the rest cannot, at least not one call at a time.

\subsection{Batch proofs over MPI}

Table~\ref{tab:mpicgbn} measures the end result for \code{gmpecpp-mpi}: the eighteen primes of
$90$ to $180$ digits in \code{bench/primes.txt}, proved with and without \code{-C}, in several
layouts of worker ranks and ECM threads, for five seeds each. Runs with and without the GPU
alternate, so slow drift affects both alike, and every proof in every run completed.

\begin{table}[t]
\centering
\begin{tabular}{lrrrr}
\toprule
Workers $\times$ threads & Seeds & Host only & \code{-C} (CGBN) & Ratio \\
\midrule
$1 \times 8$  & 5 & 47.5\,s (39.3--57.3) & 25.9\,s (23.7--27.7) & $1.83\times$ \\
$1 \times 16$ & 1 & 33.1\,s & 24.2\,s & $1.37\times$ \\
$1 \times 32$ & 1 & 27.6\,s & 23.1\,s & $1.20\times$ \\
$4 \times 8$  & 5 & 17.0\,s (13.4--19.8) & 17.4\,s (15.7--18.6) & $0.98\times$ \\
$8 \times 4$  & 5 & 21.5\,s (20.1--24.9) & 18.5\,s (17.0--19.5) & $1.16\times$ \\
\midrule
$1 \times 8$, rank bound to one core & 1 & 185.3\,s & 51.7\,s & $3.59\times$ \\
\bottomrule
\end{tabular}
\caption{Wall time of the batch of eighteen proofs, mean (minimum--maximum) over the seeds,
  with every rank unbound, except in the last row. One GPU serves every rank.}
\label{tab:mpicgbn}
\end{table}

Three things follow. First, the GPU's advantage shrinks as the host gets threads: $1.83\times$
with one worker of eight threads, $1.37\times$ with sixteen, $1.20\times$ with thirty-two. The
prover always tries its host ECM first and reaches the GPU only for cofactors of at least $192$
bits, so the more threads the host attempt has, the more often it succeeds before the GPU is
needed. Second, with four or eight workers sharing one GPU, the advantage disappears in wall
time ($0.98\times$ and $1.16\times$): the ranks contend for the device, and each pays its own
CUDA start-up. A 90-digit proof that takes $1.4$ to $1.6$\,s with \code{-C} on one worker takes
$3.2$ to $3.6$\,s with four workers and $6.5$ to $7.2$\,s with eight, across all five seeds. Third, the GPU's real contribution is predictability. Across seeds, a host-only
batch of one worker took $39$ to $57$\,s, the CGBN batch $24$ to $28$\,s.

Per proof (one worker, eight threads, median over five seeds), the picture sharpens. The GPU
turns the proofs that would grind on a stubborn cofactor into routine ones: the 140-digit proof
from $10.6$ to $1.4$\,s ($7.7\times$), a 90-digit one from $2.8$ to $0.4$\,s ($7.5\times$), the
150-digit one $5.2\times$ faster. On easy proofs it loses, by $0.47$ to $0.92\times$, to its
fixed cost per escalation; eight of the eighteen are such cases, the 180-digit proof among them.

The last row of Table~\ref{tab:mpicgbn} records a trap worth knowing. For a job of two
processes, Open MPI binds each rank to one core unless told otherwise (\code{-{}-bind-to none},
or oversubscription, which disables binding); the eight ECM threads of the worker then share
that core's two hardware threads. The host-only batch slows from $45$ to $185$\,s, while the
CGBN batch, which needs little host time, slows only from $26$ to $52$\,s, and the comparison
shows a spurious $3.6\times$. \code{mpirun -{}-report-bindings} shows the binding. The harness
passes \code{-{}-bind-to none} everywhere else.

Two properties of the prover govern how these runs should be read. Every proof reseeds the
prover's generator from the configured seed, so the order in which ranks take numbers does not
matter, and a seed fixes each proof's random choices. With more than one ECM thread, however,
the host ECM is a race --- the first thread to find a factor wins, through an atomic flag --- and
a different factor leads the downrun elsewhere; under contention from other ranks the race can
resolve differently from one run to the next. The ranges in Table~\ref{tab:mpicgbn} therefore
include variation from thread timing as well as from the seeds.

In short: CGBN helps where the work batches --- the elliptic curve method, where it is already
used --- and nowhere else. On this machine, for batch proving, it helps most when the host has
few threads per proof and the GPU few users, and it removes the occasional pathological proof in
every configuration. The $3.48\times$ of Section~\ref{sec:results} was measured with one host
thread per proof, the configuration most favourable to the GPU, and is consistent with this
trend. Moving the remaining host GMP work to the GPU would make it slower.

\section{Verification}
\label{sec:verify}

Byte-comparison against the historical program is a sound correctness
criterion only while the GPU is not engaged. Above the escalation
threshold the GPU discovers different factorisations and therefore
produces different, equally valid certificates. Two tests are accordingly
required.

Below the threshold we require byte-identity, and obtain it: 156 runs
across thirteen flag combinations, four inputs and three seeds agree on
standard output, standard error and exit status, with the backend both
enabled and disabled.

Above it we require validity, and check it with a verifier that shares no
code with the prover, its elliptic curve arithmetic written directly
against GMP. For each step it re-derives chain linkage, $q \mid m$, the
Goldwasser--Kilian bound, Hasse's bound on $m$, membership of $P$ on the
curve, equality of $[m/q]P$ with the recorded point and its distinctness
from $\mathcal{O}$, and that $[q]$ of that point is $\mathcal{O}$;
finally it requires the chain to terminate at a $q$ small enough to
settle by trial division. A self-test presents the verifier with a
certificate checked against the wrong starting value and requires
rejection, on the principle that a verifier never observed to reject
anything is not evidence.

Twenty certificates, from ten primes of ninety to one hundred and forty
digits under both backends, verify successfully.

\section{Results}
\label{sec:results}

Measurements were taken in November 2026 on an NVIDIA RTX PRO 4500 Blackwell
(sm\_120), CUDA 13.2, GCC 14.3.1 at \code{-O2}, with the seed pinned and stage two in
the prover's GPU path. The first edition's figures (RTX 4080 laptop, stage one only:
$28.8\times$ at 140 digits, the sweep $2.67\times$ in aggregate) are not directly
comparable, since the host CPUs differ (the 140-digit proof took 53\,s on that host and
68\,s on this one); like for like on this machine, the sweep's thirty-six GPU proofs took
$145.9$\,s with stage one only and $117.9$--$120.8$\,s with stage two.

\begin{table}[t]
\centering
\begin{tabular}{rrrrrr}
\toprule
Digits & Host & GPU & Speedup & Host ECM & GPU ECM \\
\midrule
20  & 0.008\,s  & 0.008\,s & 1.01$\times$ & 0.004\,s & 0.004\,s \\
30  & 0.110\,s  & 0.108\,s & 1.02$\times$ & 0.064\,s & 0.065\,s \\
40  & 0.173\,s  & 0.175\,s & 0.99$\times$ & 0.156\,s & 0.158\,s \\
60  & 0.209\,s  & 0.209\,s & 1.00$\times$ & 0.159\,s & 0.159\,s \\
90  & 2.698\,s  & 1.067\,s & 2.53$\times$ & 2.580\,s & 0.973\,s \\
110 & 6.574\,s  & 2.009\,s & 3.27$\times$ & 6.373\,s & 1.856\,s \\
140 & 67.994\,s & 2.660\,s & 25.57$\times$ & 67.787\,s & 2.444\,s \\
\bottomrule
\end{tabular}
\caption{Whole-proof timings. The final two columns confirm that the
  entire difference is attributable to the factoring stage.}
\label{tab:proofs}
\end{table}

Over a thirty-six-proof sweep covering every prime of ninety digits or
more in the committed benchmark data, at two seeds each, the GPU backend
is faster in twenty-seven runs, with median speedup $2.37\times$, range
$0.57\times$ to $22.6\times$, and aggregate $420.6 \to 120.8$ seconds,
or $3.48\times$.

The aggregate understates the phenomenon. Host runtimes across that sweep
span $0.83$ to $67.3$ seconds; GPU-assisted runtimes span $0.91$ to
$19.7$ seconds. The effect is best described as variance reduction: the backend
does not uniformly accelerate proofs so much as eliminate pathological
ones, in which the host's twenty curves per bound labour against a
cofactor that thousands of curves split at once. Every run in which the
backend loses is one in which the host took $6.5$\,s or less, except one 180-digit
run ($19.3$ against $19.7$\,s); at 180 digits the two are level.

A corollary appears in the certificate lengths. At one hundred and forty
digits the host required fifteen steps and $51.9$ seconds where the GPU
found a superior factorisation in eight steps and $1.9$ seconds: a
discriminant abandoned by the host for want of a factorisation is one the
GPU can use, so its proofs are frequently shorter as well as faster.

\section{Limitations}

Class polynomial degree, that is the class number $h(D)$, must remain
below $8192$; no bounds check exists. Polynomial buffers are fixed-size
stack arrays and the root-finding routine is recursive, at approximately
$512$\,KB per frame. The CGBN stage two uses the simple prime
continuation, whose cost grows with the number of primes up to $B_2$; GMP-ECM's FFT
continuation reaches far larger $B_2$ for the same work. The cuFFT backend has that
continuation (Section~\ref{sec:fft}) but only above $16352$ bits, and its fixed costs remain
large: every coefficient of every tree node costs a Montgomery reduction. Batching each tree
level into one transform, and computing baby and giant points in batched rounds rather than one
serial chain, removed per-operation launch costs, which matter only with few curves.
The exception introduced in
place of \code{exit(2)} unwinds frames holding raw GMP temporaries, which
leak on that path; the path is reached only for composite inputs, but the
defect should be addressed before the library is applied to many
composites in one process.

The serialisation of API entry points remains, so a single process cannot
have two proofs in flight; within one proof it does use many cores, but
callers wanting concurrent proofs need separate processes or the MPI
driver. Removing that restriction means threading a context through the
file-scope globals, which is the same refactor the prover has resisted
since it was a single translation unit.

Two parallel axes identified in Table~\ref{tab:axes} are unexploited.
The $m$ candidates for a given discriminant, of which there are between
two and six, are independent factoring attempts; the fan-out is small and
the sequential form short-circuits, so the gain is bounded. More
promising is the class polynomial layer, which is a map over reduced
forms followed by a reduction, with each $j(\tau)$ evaluation expensive
and independent. It is worth doing at high precision, where that layer
dominates; at the working default it is a tenth of a percent to a third
of runtime and the factoring stage is the better target.

\section*{Acknowledgements}

The prover is the work of Pate Williams, James Wanless, Chris Studholme,
Elisavet Konstantinou, Yiannis Stamatiu, Christos Zaroliagis, Scott
Contini, Paul Zimmermann and Timo Poikola, to whom the substance of this
software is owed. We thank the authors of GMP and of CGBN.

\begin{thebibliography}{99}

\bibitem{atkin1993} A. O. L. Atkin and F. Morain.
  \emph{Elliptic curves and primality proving}.
  Mathematics of Computation, 61(203):29--68, 1993.

\bibitem{bernstein2004} D. J. Bernstein.
  \emph{Scaled remainder trees}. Preprint, 2004.
  \url{https://cr.yp.to/papers.html#scaledmod}.

\bibitem{cantor1981} D. G. Cantor and H. Zassenhaus.
  \emph{A new algorithm for factoring polynomials over finite fields}.
  Mathematics of Computation, 36(154):587--592, 1981.

\bibitem{cgbn} NVIDIA Labs.
  \emph{CGBN: CUDA Accelerated Multiple Precision Arithmetic (Big Num)
  using Cooperative Groups}.
  \url{https://github.com/NVlabs/CGBN}.

\bibitem{cohen1993} H. Cohen.
  \emph{A Course in Computational Algebraic Number Theory}.
  Graduate Texts in Mathematics 138. Springer-Verlag, 1993.

\bibitem{cornacchia1908} G. Cornacchia.
  \emph{Su di un metodo per la risoluzione in numeri interi
  dell'equazione $\sum_{h=0}^{n} C_h x^{n-h} y^h = P$}.
  Giornale di Matematiche di Battaglini, 46:33--90, 1908.

\bibitem{cufft} NVIDIA Corporation.
  \emph{cuFFT Library User's Guide}, CUDA Toolkit 13.2.
  \url{https://docs.nvidia.com/cuda/cufft/}.

\bibitem{gmp} T. Granlund and the GMP development team.
  \emph{GNU MP: The GNU Multiple Precision Arithmetic Library}.
  \url{https://gmplib.org/}.

\bibitem{goldwasser1986} S. Goldwasser and J. Kilian.
  \emph{Almost all primes can be quickly certified}.
  In Proceedings of the 18th Annual ACM Symposium on Theory of
  Computing, pages 316--329, 1986.

\bibitem{ieee1363} IEEE.
  \emph{IEEE Standard Specifications for Public-Key Cryptography}.
  IEEE Std 1363-2000, 2000.

\bibitem{kaltofen1989} E. Kaltofen and T. Valente.
  \emph{An improved Las Vegas primality test}.
  In Proceedings of the ACM-SIGSAM International Symposium on Symbolic
  and Algebraic Computation, pages 26--33, 1989.

\bibitem{kaltofen1991} E. Kaltofen and N. Yui.
  \emph{Explicit construction of the Hilbert class fields of imaginary
  quadratic fields by integer lattice reduction}.
  In Number Theory (New York Seminar 1989--1990), pages 149--202.
  Springer-Verlag, 1991.

\bibitem{konstantinou2002} E. Konstantinou, Y. C. Stamatiou and
  C. Zaroliagis.
  \emph{On the efficient generation of elliptic curves over prime
  fields}.
  In Cryptographic Hardware and Embedded Systems (CHES 2002),
  Lecture Notes in Computer Science 2523, pages 333--348, 2003.

\bibitem{lenstra1987} H. W. Lenstra, Jr.
  \emph{Factoring integers with elliptic curves}.
  Annals of Mathematics, 126(3):649--673, 1987.

\bibitem{miller1976} G. L. Miller.
  \emph{Riemann's hypothesis and tests for primality}.
  Journal of Computer and System Sciences, 13(3):300--317, 1976.

\bibitem{montgomery1985} P. L. Montgomery.
  \emph{Modular multiplication without trial division}.
  Mathematics of Computation, 44(170):519--521, 1985.

\bibitem{montgomery1987} P. L. Montgomery.
  \emph{Speeding the Pollard and elliptic curve methods of
  factorization}.
  Mathematics of Computation, 48(177):243--264, 1987.

\bibitem{montgomery1992} P. L. Montgomery.
  \emph{An FFT extension of the elliptic curve method of factorization}.
  PhD thesis, University of California, Los Angeles, 1992.

\bibitem{niven1991} I. Niven, H. S. Zuckerman and H. L. Montgomery.
  \emph{An Introduction to the Theory of Numbers}, 5th edition.
  John Wiley and Sons, 1991.

\bibitem{rabin1980} M. O. Rabin.
  \emph{Probabilistic algorithm for testing primality}.
  Journal of Number Theory, 12(1):128--138, 1980.

\bibitem{vzgathen2013} J. von zur Gathen and J. Gerhard.
  \emph{Modern Computer Algebra}, 3rd edition.
  Cambridge University Press, 2013.

\bibitem{weber1908} H. Weber.
  \emph{Lehrbuch der Algebra}, Band III, 2nd edition.
  Friedrich Vieweg und Sohn, 1908.

\bibitem{zimmermann2006} P. Zimmermann and B. Dodson.
  \emph{20 years of ECM}.
  In Algorithmic Number Theory (ANTS-VII), Lecture Notes in Computer
  Science 4076, pages 525--542, 2006.

\end{thebibliography}

\end{document}
